% I. Анализ на предметната област и литературен преглед.
% Покрива темите от учебната програма: характеристики и особености на естествените
% езици, математически и психологически модели, модели за представяне на знания.
\section{Анализ на предметната област и литературен преглед}
\label{sec:literature}

\subsection{Характеристики и особености на естествените езици}

Естественият език се различава от формалните езици по няколко свойства, всяко от
които е пряка причина за инженерна трудност. Първото е \term{нееднозначността}: една
и съща поредица от думи допуска повече от едно разчитане на всяко ниво, лексикално,
морфологично, синтактично и семантично. Второто е \term{продуктивността}: множеството
от допустими изречения е безкрайно, така че системата не може да работи със списък от
разпознати въпроси. Третото е \term{контекстната зависимост}: значението на дадено
изречение се променя от предходния текст и от предметната област.

Към тези свойства се добавят особеностите на българския език, който е силно
флективен. Една лексема поражда десетки словоформи, а членуването, родът и числото
носят синтактична информация, която в английския се изразява с ред на думите или със
служебни думи. Това измества тежестта на анализа към морфологичното ниво: без
надежден морфологичен анализатор синтактичният анализатор получава вход, в който
почти всяка дума е неразпозната.

\subsection{Математически модели на естествения език}

Математическото описание на езика започва от йерархията на формалните граматики и от
въпроса къде в нея попада естественият език. Класическата постановка е дадена
в~\cite{chomsky1957syntactic} и въвежда генеративната граматика като модел на
езиковата компетентност. За целите на инженерната реализация практическият избор
почти винаги е безконтекстна граматика, разширена с признаци, защото за нея са
известни ефективни алгоритми за анализ~\cite{earley1970parsing,younger1967cyk}, а
изразителността ѝ покрива преобладаващата част от въпросителните конструкции.

На морфологично ниво съответният математически апарат е този на крайните автомати и
на крайните преобразуватели. Преобразувателят изобразява повърхнинна словоформа върху
лексикална форма заедно с граматични признаци и обратно, като едно и също устройство
служи както за анализ, така и за синтез~\cite{mohri1997fst,beesley2003fsmorphology}.
Това двупосочно свойство е причината същият апарат да покрие и задачата за синтез на
изречение в Раздел~\ref{sec:design}.

\subsection{Психологически модели на езиковата обработка}

Психолингвистичните модели описват как човек обработва изречение в реално време и
имат пряко отражение върху избора на алгоритъм. Инкременталността, тоест изграждането
на частична интерпретация още преди изречението да е завършило, е наблюдавана при
човешкия анализ и е свойство и на алгоритъма на Earley~\cite{earley1970parsing}.
Мрежите с преходи~\cite{woods1970atn} са ранен опит да се съчетае изчислителен модел с
психологическа правдоподобност. \TODO{допълнителен източник за психологическите модели
на разбиране на изречение, сверен по издателска страница}

\subsection{Модели за представяне на знания и съпоставителна оценка}

Формалната заявка се произвежда срещу някакъв модел на знанието, така че изборът на
този модел определя какво изобщо може да бъде попитано. Основните семейства са
логиката от първи ред, семантичните мрежи и концептуалните графи, фреймовите
представяния, релационният модел и графовите модели от типа
RDF~\cite{sowa2000kr,w3c2013sparql}. Те се различават по изразителност, по цената на
извода и по това колко естествено върху тях се изобразява резултатът от семантичния
анализ. Съпоставителната оценка по тези критерии е дадена в
Раздел~\ref{sec:alternatives}, където се обосновава и направеният избор.

\subsection{Интерфейси на естествен език към бази от данни}

Задачата има дълга история като самостоятелно направление. Обзорът
в~\cite{androutsopoulos1995nlidb} систематизира архитектурите на символните системи и
описва повтарящите се проблеми, преносимост към нова предметна област, обработка на
нееднозначност и обратна връзка към потребителя. По-новите набори от данни и мерки за
оценка~\cite{zelle1996geoquery,yu2018spider} дават общ терен, върху който символни и
статистически решения могат да се сравняват. Обобщение на прегледаните подходи и
позиционирането на настоящата разработка спрямо тях е дадено в Табл.~\ref{tab:systems}.

\begin{table}[H]
\caption{Прегледани подходи и позициониране на настоящата разработка}
\label{tab:systems}
\centering
\begin{tabular}{@{}L{3.2cm}L{2.6cm}L{3.4cm}L{3.6cm}@{}}
\toprule
\textbf{Подход} & \textbf{Метод} & \textbf{Нееднозначност} & \textbf{Обяснимост} \\
\midrule
мрежи с преходи~\cite{woods1970atn} & символен, процедурен & чрез връщане назад в мрежата & следата е мрежата, четима, но процедурна \\
обзор на символните интерфейси~\cite{androutsopoulos1995nlidb} & символен, слоест & разрешава се на семантично ниво & по построение \\
индуктивно логическо програмиране~\cite{zelle1996geoquery} & обучение върху двойки & статистическа & частична, чрез научените правила \\
съвременни набори за текст към SQL~\cite{yu2018spider} & статистически & неявна & липсва проверима следа \\
настоящата разработка & символен, слоест & уплътнена гора и типова проверка & следа от думата до заявката, с причина за отказа \\
\bottomrule
\end{tabular}
\end{table}

Отделно като свързана работа стои изследване на автора върху съпоставянето с
регулярни изрази, което ползва същия апарат от крайни автомати като морфологичния
слой тук. \TODO{цитат на статията за съпоставяне с регулярни изрази}
